\documentclass[a4paper,12pt]{article}
\setcounter{secnumdepth}{5}
\setcounter{tocdepth}{3}
\input{/usr/share/LaTeX-ToolKit/template.tex}
\begin{document}
\title{晶體}
\author{沈威宇}
\date{\temtoday}
\titletocdoc
\sct{晶體（Crystal）}
成分粒子（如原子、分子或離子）以高度有序的微觀結構排列的固體，形成向各個方向延伸的晶格。 
\ssc{晶胞（Cell）/晶格（Lattice）}
指構成晶體重複結構的單位。
\sssc{原始晶胞（primitive cell）}
指構成晶體重複結構的最小單位。其表示方法不指一種，但均包含相同數量的晶格點，且對於二維晶體，其面積固定，為四個頂點共有一個粒子的平行四邊形；對於三維晶體，其體積固定，為八個頂點共有一個粒子的平行六面體。
\sssc{配位數}
與一個晶格點最近的晶格點數量。
\sssc{維格納-賽茨晶胞（Wigner–Seitz cell）}
對於二維晶體，一晶格點（lattice point）周圍的維格納-賽茨晶胞被定義為平面上距離該晶格點比距離任何其他晶格點更近的點的軌跡。

對於三維晶體，一晶格點周圍的維格納-賽茨晶胞被定義為空間中距離該晶格點比距離任何其他晶格點更近的點的軌跡。
\sssc{原始平移向量（Primitive translation vectors）}
二維晶體晶格的原始平移向量為$\mathbf {a}$與$\mathbf{b}$表示以晶體中一原子為原點，晶體中的原子均位於$p\mathbf {a} +q\mathbf {b} $，其中$p$、 $q$為整數。

三維晶體晶格的原始平移向量為$\mathbf {a}$、$\mathbf{b}$與$\mathbf {c}$表示以晶體中一原子為原點，晶體中的原子均位於$p\mathbf {a} +q\mathbf {b} +r\mathbf {c} $，其中$p$、 $q$、$r$為整數。
\ssc{晶族（crystal family）、晶系（crystal system）與晶格系（lattice system）}
\sssc{分類}
\begin{longtable}[c]{|p{0.2\tw}|p{0.2\tw}|p{0.2\tw}|p{0.2\tw}|}
    \hline
        晶族 & 晶系（由點群決定） & 點群所需的對稱性 & 晶格系（由布拉維格子決定） \\ \hline\endhead
        三斜（Triclinic） & 三斜（Triclinic） & 無 & 三斜（Triclinic） \\ \hline
        單斜（Monoclinic） & 單斜（Monoclinic） & 1個雙重旋轉軸或1個鏡像平面 & 單斜（Monoclinic） \\ \hline
        斜方/正交（Orthorhombic） & 斜方/正交（Orthorhombic） & 3個雙重旋轉軸或1個雙重旋轉軸和2個鏡像平面 & 斜方/正交（Orthorhombic） \\ \hline
        四方/正方（Tetragonal） & 四方/正方（Tetragonal） & 1個四重旋轉軸 & 四方/正方（Tetragonal） \\ \hline
        六方（Hexagonal） & 三方（Trigonal） & 1個三重旋轉軸 & 菱面（Rhombohedral）或六方（Hexagonal） \\ \hline
        六方（Hexagonal） & 六方（Hexagonal） & 1個六重旋轉軸 & 六方（Hexagonal） \\ \hline
        立方/等軸（Cubic） & 立方/等軸（Cubic） & 4個三重旋轉軸 & 立方/等軸（Cubic） \\ \hline
    \end{longtable}\FB
\sssc{幾何詞彙釋義}
\bit
\item \tb{三維最密堆積/密排（close-packed）}：最密等體積球填充，堆積密度 $\frac{\sqrt{2}\pi}{6}\approx 74\%$。
\item \tb{二維最密堆積/密排（close-packed）}：最密等面積圓填充，堆積密度 $\frac{\sqrt{3}\pi}{6}\approx 91\%$。
\item \tb{截角八面體（Truncated octahedron）}：正八面體的六個頂點截去六個四角錐形成的具有14個面（8個全等正六邊形與6個全等正方形）、36條邊與24個頂點的凸多面體。
\item \tb{菱形十二面體（Rhombic dodecahedron）}：具有12個全等菱形面、24條邊與14個頂點的凸多面體。
\eit
\sssc{單成分粒子晶體}
\begin{longtable}[c]{|p{0.15\textwidth}|p{0.05\textwidth}|p{0.2\textwidth}|p{0.05\textwidth}|p{0.05\textwidth}|p{0.1\textwidth}|p{0.1\textwidth}|p{0.1\textwidth}|}
\hline
堆積方式 & 配位數 & 平行四邊形/平行六面體晶胞描述 & 平行四邊形/平行六面體晶胞粒子數 & 平行四邊形/平行六面體晶胞邊長除以粒子半徑 & 維格納-賽茨晶胞形狀 & 堆積密度 & 舉例 \\\hline\endhead
\tb{簡單立方（Simple cubic）/原始立方（Primitive cubic, cP）} & $6$ & 每個頂點各八分之一個粒子 & $1$ & $2$ & 正方體 & $\frac{\pi}{6}\approx 52\%$ & 釙（唯一已知自然以簡單立方存在之金屬） \\\hline
\tb{面心立方（Face-centered cubic, cF, fcc）/立方密排（cubic closepacked）} & $12$ & 簡單立方每面中心再加半個粒子 & $4$ & $2\sqrt{2}$ & 截角八面體 & $\frac{\sqrt{2}\pi}{6}\approx 74\%$ & IA、鈣、鍶、鋁、鎳、銀、銅、金 \\\hline
\tb{體心立方（Body-centered cubic, cL, bcc）} & $8$ & 簡單立方中心再加一個粒子 & $2$ & $\frac{4\sqrt{3}}{3}$ & 菱形十二面體 & $\frac{\sqrt{3}\pi}{8}\approx 68\%$ & 鋇、鐳、鐵、釩、鈮、鉻 \\\hline
\tb{鑽石立方（diamond cubic, dc）} & $4$ & 面心立方，取四個兩兩為一面對角線的頂點，分別與與該頂點相鄰的三個面的中心做一四面體，在該四個四面體的中心各增加一個粒子 & $8$ & $\frac{8\sqrt{3}}{3}$ & - & $\frac{\sqrt{3}\pi}{16}\approx 34\%$ & 金剛石、矽、碳化矽、α-錫（灰錫） \\\hline
\tb{簡單六方（Simple hexagonal）} & $8$ & 角度60°與120°的菱形底面的菱形柱，底面60°的頂點各 $\frac{1}{12}$ 個粒子、底面120°的頂點各 $\frac{1}{6}$ 個粒子 & $1$ & $2$ & 正六邊形柱 & $\frac{\sqrt{3}\pi}{9}\approx 61\%$ & - \\\hline
\tb{六方密排/六方最密（Hexagonal close-packed, hcp）} & $12$ & 簡單六方，其菱形柱分成兩個正三角柱，其一的中心加一粒子 & $2$ & 菱形底面$2$、高$\frac{8\sqrt{3}}{3}$ & - & $\frac{\sqrt{2}\pi}{6}\approx 74\%$ & 鈹、鎂、鈦、鈷、鋅、鎘、鈧、釔 \\\hline
\end{longtable}\FB
\sssc{一比一雙成分粒子晶體}
\begin{longtable}[c]{|p{0.15\textwidth}|p{0.05\textwidth}|p{0.2\textwidth}|p{0.05\textwidth}|p{0.05\textwidth}|p{0.1\textwidth}|p{0.1\textwidth}|p{0.1\textwidth}|}
\hline
堆積方式 & 配位數 & 平行四邊形/平行六面體晶胞描述 & 平行四邊形/平行六面體晶胞粒子數 & 平行四邊形/平行六面體晶胞邊長除以兩成分粒子半徑和 & 維格納-賽茨晶胞形狀 & 堆積密度 & 舉例 \\\hline\endhead
\tb{平面三角（Trigonal planar）} & $3$ & 角度60°與120°的菱形，60°的頂點各 $\frac{1}{6}$ 個一類粒子、120°的頂點各 $\frac{1}{3}$ 個該類粒子、一60°的頂點與其二相鄰頂點形成的正三角形的中心一個另一類粒子 & $2$ & $\sqrt{3}$ & 正三角形 & $\frac{\sqrt{3}\pi}{9}\approx 61\%$ & - \\\hline
\tb{B1/NaCl/氯化鈉/岩鹽} & $6$ & 一類粒子面心立方、另一類粒子置於邊中點與體心，屬互穿面心立方 & $8$ & $2$ & 正方體 & $\frac{\pi}{6}\approx 52\%$ & \ce{NaCl}, \ce{MgO} \\\hline
\tb{B2/CsCl/氯化銫} & $8$ & 一類粒子簡單立方、另一類粒子置於體心，屬互穿簡單立方 & $2$ & $\frac{2\sqrt{3}}{3}$ & 正八面體 & $\frac{\sqrt{3}\pi}{8}\approx 68\%$ & \ce{CsCl}, \ce{NH4Br} \\\hline
\tb{B3/ZnS/硫化鋅/閃鋅礦} & $4$ & 一類粒子面心立方，取四個兩兩為一面對角線的頂點，分別與與該頂點相鄰的三個面的中心做一四面體，在該四個四面體的中心各增加一個另一類粒子，屬互穿面心立方 & $8$ & $\frac{4\sqrt{3}}{3}$ & 正四面體 & $\frac{\sqrt{3}\pi}{16}\approx 34\%$ & \ce{ZnS}, \ce{CuCl} \\\hline
\end{longtable}\FB
\end{document}
